% !TEX root = ../main.tex

\ifprmresultsready

Impurity formation-energy transfer within IMP2D is organized by chemical
novelty.  Pair-held-out folds, with one singleton-host exception, gave
\PRMPairCvDARTMAE\eV{} MAE, compared with \PRMIdCvDARTMAE\eV{} under random
cross-validation.  Removing all target supervision for a host or impurity
increased the error to \PRMHostCvDARTMAE\eV{} and
\PRMDopantCvDARTMAE\eV{}, respectively.  Incorporation geometry added a second
axis: interstitial structures were substantially harder to predict than
adsorbates in the pair-held-out population.

The architecture results connect graph pooling to a formation-energy target
that is non-extensive in host atom count.  Gated graph readout and the composite
pre-normalized, distance-gated local block produced small, repeat-consistent
gains, whereas explicit environment enrichment had no resolved independent
effect.  In the evaluated SchNet recipe, a post hoc sensitivity test implicated
atomwise addition, which is natural for an extensive total energy, as one
contributor to the large host-held-out error.

For existing DFT-relaxed candidate geometries, pair-held-out predictions
selected the lower-energy incorporation class with \PRMPreferenceAccuracy\%
accuracy and \PRMGlobalRegret\eV{} mean global site-selection regret.  Because
the geometries and reference energies arise from the same completed DFT
workflow, these results support retrospective energy estimation and reranking
within the neutral, retained IMP2D domain rather than a reduction in relaxation
cost.  Testing newly generated structures prepared independently of their
reference energies, together with uncertainty calibration under explicit
chemical shift, remains necessary before prospective use beyond that domain.

\else

\noindent\textit{The numerical conclusion is unavailable because the audited
result bundle is not present.}

\fi
